\chapter{System Operation Flow\label{chap:flow}}

This chapter describes the complete operation flow of the Catch Turtle All system, from startup to continuous operation.

\section{Operation Phases}

The system operates in a continuous loop through the following phases:

\begin{enumerate}[label=\textbf{Phase \arabic*:}]
  \item \textbf{Startup}: The launch file starts all nodes simultaneously. \texttt{turtlesim\_node} creates \texttt{turtle1} at the centre of the canvas (coordinates 5.544455, 5.544455). The \texttt{spawn\_manager} waits for the \texttt{/spawn} service to become available before starting its timer.

  \item \textbf{Target Generation}: \texttt{spawn\_manager} begins spawning new turtles every 3 seconds at random positions within the configured bounds. Each new turtle immediately starts publishing its pose on the \texttt{/turtleN/pose} topic.

  \item \textbf{Target Discovery}: \texttt{master\_manager} periodically scans the ROS2 topic list (every 1.0\,s) and subscribes to any new \texttt{/turtleN/pose} topics it discovers. It updates the \texttt{TurtleRegistry} with each turtle's position as pose messages arrive.

  \item \textbf{Target Selection}: Every 0.5 seconds, \texttt{master\_manager} evaluates all uncaught turtles---excluding the master turtle, already-caught turtles, and cooled-down targets---and selects the one nearest to the master turtle using Euclidean distance.

  \item \textbf{Pursuit Execution}: \texttt{master\_manager} sends a \texttt{CatchTarget} goal to \texttt{catch\_executor}, which takes control of \texttt{turtle1}'s velocity and drives it toward the target. Real-time distance feedback is published back to \texttt{master\_manager} at 20\,Hz.

  \item \textbf{Preemption Check}: During pursuit, \texttt{master\_manager} continuously re-evaluates target proximity. If a closer target appears (satisfying the hysteresis condition $d_{\text{new}} + m_{\text{preempt}} < d_{\text{current}}$), the current goal is cancelled and a new one is dispatched for the closer target on the next decision tick.

  \item \textbf{Catch Completion}: When \texttt{catch\_executor} detects that the distance to the target is below the catch threshold (0.5\,m), it marks the goal as succeeded, publishes a zero-velocity command to stop the master turtle, and returns the result.

  \item \textbf{Chain Update}: \texttt{master\_manager} receives the successful result, marks the caught turtle in the registry, appends it to the chain, and publishes the updated chain on \texttt{/caught\_chain} as a JSON message.

  \item \textbf{Formation Update}: \texttt{follower\_manager} receives the chain update, creates the necessary pose subscriber and velocity publisher for the new follower, and begins controlling it to follow its predecessor at the configured following distance.

  \item \textbf{Loop}: The system returns to Phase~4 and continues the pursuit cycle indefinitely.
\end{enumerate}

\section{Edge Cases and Error Handling}

The system handles several edge cases gracefully:

\begin{itemize}
  \item \textbf{No uncaught targets}: If no uncaught turtles are available, \texttt{master\_manager} simply waits and re-evaluates on the next decision tick. No action goal is dispatched.

  \item \textbf{Action server unavailable}: If \texttt{catch\_executor} is not yet ready, \texttt{master\_manager} logs a warning and retries on the next decision tick.

  \item \textbf{Goal rejection}: If \texttt{catch\_executor} rejects a goal (because it is busy), \texttt{master\_manager} places the target on cooldown and moves on to the next candidate.

  \item \textbf{Goal timeout}: If a pursuit takes longer than \texttt{goal\_timeout\_sec}, \texttt{catch\_executor} aborts the goal and \texttt{master\_manager} places the target on cooldown.

  \item \textbf{Missing pose data}: If no pose is received for either the master or target turtle within \texttt{no\_pose\_timeout\_sec}, the goal is aborted.

  \item \textbf{Spawn failure}: If \texttt{turtlesim} rejects a spawn request, \texttt{spawn\_manager} increments the index and retries on the next tick.
\end{itemize}
